% ---- EDy4: Stationäres Strömungsfeld --------------------

% 4.1  Stromdichte
\bereich[cKap4]{4.1 Stromdichte $\vec{J}$}
\begin{formelreihe}
  \parbox[t]{\linewidth}{%
    \formelblock{Stromdichte (allgemein)}{%
      I=\int_A\!\vec{J}\,\mathrm{d}\vec{A}}\\[0.25em]%
    \formelblock{Homogenes Feld}{%
      J=\tfrac{I}{A}\;,\qquad [J]=1\,\tfrac{\mathrm{A}}{\mathrm{m}^2}}} &
  \parbox[t]{\linewidth}{%
    \formelblock{Stromdichtevektor}{%
      \vec{J}\parallel\vec{E}\;,\quad
      \vec{J}\parallel\vec{S}_{\mathrm{Feldl.}}}\\[0.25em]%
    \infoblock{Strömungsfeld quellenfrei:\\
      $\oint\vec{J}\,\mathrm{d}\vec{A}=0$}} &
  \beispielblock{Draht, $d=2\,\text{mm}$, $I=1\,\text{A}$}{%
    \begin{aligned}
      A&=\pi r^2=\pi(1\,\text{mm})^2\\[0.1em]
      J&=\tfrac{1\,\text{A}}{\pi\cdot10^{-6}\,\text{m}^2}\\[0.1em]
      &\approx318\,\tfrac{\text{kA}}{\text{m}^2}
    \end{aligned}} \\[0.3em]
\end{formelreihe}
\bereichende

% 4.2  Elektrische Leitfähigkeit / Lok. Ohmsches Gesetz
\bereich[cKap4]{4.2 Elektrische Leitfähigkeit — Lokales Ohmsches Gesetz}
\begin{formelreihe}
  \parbox[t]{\linewidth}{%
    \formelblock{Lokales Ohmsches Gesetz}{%
      \vec{J}=\kappa\cdot\vec{E}}\\[0.25em]%
    \formelblock{Spez.\ Leitfähigkeit}{%
      [\kappa]=1\,\tfrac{\mathrm{S}}{\mathrm{m}}
        =1\,\tfrac{\mathrm{A}}{\mathrm{V\,m}}}} &
  \parbox[t]{\linewidth}{%
    \formelblock{Spez.\ Widerstand}{%
      \varrho=\tfrac{1}{\kappa}\;,\qquad
      [\varrho]=1\,\Omega\mathrm{m}}\\[0.25em]%
    \infoblock{Kupfer: $\kappa_\mathrm{Cu}=58\cdot10^6\,\tfrac{\mathrm{S}}{\mathrm{m}}$\\
      Isolator: $\kappa\to0$, Metall: $\kappa\gg1$}} &
  \beispielblock{Analogie}{%
    \begin{aligned}
      \vec{J}&=\kappa\,\vec{E}\\[0.1em]
      \vec{D}&=\varepsilon\,\vec{E}\\[0.1em]
      \Rightarrow\;\kappa&\leftrightarrow\varepsilon
    \end{aligned}} \\[0.3em]
\end{formelreihe}
\bereichende

% 4.3  Kontinuitätsgleichung
\bereich[cKap4]{4.3 Kontinuitätsgleichung (Stationär)}
\begin{formelreihe}
  \parbox[t]{\linewidth}{%
    \formelblock{Quellfreiheit der Stromdichte}{%
      \oint_{A}\!\vec{J}\,\mathrm{d}\vec{A}=0}\\[0.25em]%
    \formelblock{Differenzielle Form}{%
      \nabla\cdot\vec{J}=0}} &
  \parbox[t]{\linewidth}{%
    \formelblock{Knotenpunktsatz (Kirchhoff)}{%
      \textstyle\sum_k I_k=0}\\[0.25em]%
    \infoblock{Stationär: keine freien Ladungen\\
      im Leiterinneren ($\varrho_\mathrm{frei}=0$)}} &
  \beispielblock{Zwei Leiter, $A_1\neq A_2$}{%
    \begin{aligned}
      I_1&=I_2=I\\[0.1em]
      J_1 A_1&=J_2 A_2\\[0.1em]
      \Rightarrow J&\sim\tfrac{1}{A}
    \end{aligned}} \\[0.3em]
\end{formelreihe}
\bereichende

% 4.4  Widerstand — Globales Ohmsches Gesetz
\bereich[cKap4]{4.4 Widerstand — Globales Ohmsches Gesetz}
\begin{formelreihe}
  \parbox[t]{\linewidth}{%
    \formelblock{Ohmsches Gesetz (integral)}{%
      U=R\cdot I\;,\qquad [R]=1\,\Omega}\\[0.25em]%
    \formelblock{Widerstand aus Geometrie}{%
      R=\varrho\cdot\tfrac{l}{A}=\tfrac{l}{\kappa\cdot A}}} &
  \parbox[t]{\linewidth}{%
    \formelblock{Leitwert}{%
      G=\tfrac{1}{R}=\kappa\cdot\tfrac{A}{l}\;,\quad[G]=1\,\mathrm{S}}\\[0.25em]%
    \formelblock{Herleitung ($\vec{J}=\kappa\vec{E}$, homogen)}{%
      U=E\cdot l\;,\;I=J\cdot A=\kappa E A}} &
  \beispielblock{Zylinder, $l$, $A$, $\kappa$}{%
    \begin{aligned}
      R&=\tfrac{l}{\kappa A}\\[0.1em]
      U&=\tfrac{I\,l}{\kappa A}\\[0.1em]
      E&=\tfrac{U}{l}=\tfrac{J}{\kappa}
    \end{aligned}} \\[0.3em]
\end{formelreihe}
\bereichende

% 4.5  Verlustleistung (Joulesche Wärme)
\bereich[cKap4]{4.5 Verlustleistung — Joulesche Wärme}
\begin{formelreihe}
  \parbox[t]{\linewidth}{%
    \formelblock{Leistung (global)}{%
      P=U\cdot I=R\cdot I^2=\tfrac{U^2}{R}}\\[0.25em]%
    \formelblock{Leistungsdichte (lokal)}{%
      p=\kappa\cdot E^2=\vec{J}\cdot\vec{E}=\tfrac{J^2}{\kappa}}} &
  \parbox[t]{\linewidth}{%
    \formelblock{Gesamtleistung aus Dichte}{%
      P=\int_V\!p\,\mathrm{d}V
        =\int_V\!\kappa E^2\,\mathrm{d}V}\\[0.25em]%
    \infoblock{Elektrische Energie → Wärme\\
      (Joulsche Verluste, irreversibel)}} &
  \beispielblock{Zylinder, $\kappa$, $E$, $V=A\cdot l$}{%
    \begin{aligned}
      p&=\kappa E^2\\[0.1em]
      P&=\kappa E^2\cdot A l\\[0.1em]
      &=\tfrac{U^2}{R}
    \end{aligned}} \\[0.3em]
\end{formelreihe}
\bereichende

% 4.6  Grenzflächen im Strömungsfeld
\bereich[cKap4]{4.6 Grenzflächen im Strömungsfeld}
\begin{formelreihe}
  \parbox[t]{\linewidth}{%
    \formelblock{Normalkomp.\ stetig ($\oint\vec{J}\,\mathrm{d}\vec{A}=0$)}{%
      J_{n1}=J_{n2}}\\[0.25em]%
    \formelblock{Tangentialkomp.\ stetig ($\oint\vec{E}\,\mathrm{d}\vec{s}=0$)}{%
      E_{t1}=E_{t2}}} &
  \parbox[t]{\linewidth}{%
    \formelblock{Brechungsgesetz}{%
      \tfrac{\tan\alpha_1}{\tan\alpha_2}=\tfrac{\kappa_1}{\kappa_2}}\\[0.25em]%
    \infoblock{Analog zum Magnetfeld:\\
      $\tfrac{\tan\alpha_1}{\tan\alpha_2}=\tfrac{\mu_1}{\mu_2}$\quad
      ($J_{n}\leftrightarrow B_{n}$)}} &
  \beispielblock{Grenzfl.\ $\kappa_1\gg\kappa_2$}{%
    \begin{aligned}
      &\kappa_1\gg\kappa_2:\\[0.1em]
      &\Rightarrow\alpha_1\gg\alpha_2\\[0.1em]
      &\Rightarrow\vec{J}\text{ senkr.\ in Med.\ 2}
    \end{aligned}} \\[0.3em]
\end{formelreihe}
\bereichende

% 4.7  Analogie Strömungsfeld ↔ Elektrostatik
\bereich[cKap4]{4.7 Analogie: Strömungsfeld $\leftrightarrow$ Elektrostatik}
\begin{formelreihe}
  \parbox[t]{\linewidth}{%
    \formelblock{Feldgrößen-Analogie}{%
      \vec{J}=\kappa\vec{E}
        \;\leftrightarrow\;
      \vec{D}=\varepsilon\vec{E}}\\[0.25em]%
    \formelblock{Quell-/Wirbelfreiheit}{%
      \oint\vec{J}\,\mathrm{d}\vec{A}=0
        \;\leftrightarrow\;
      \oint\vec{D}\,\mathrm{d}\vec{A}=0}} &
  \parbox[t]{\linewidth}{%
    \formelblock{Schaltgrößen-Analogie}{%
      G\leftrightarrow C\;,\quad R\leftrightarrow\tfrac{1}{C}}\\[0.25em]%
    \formelblock{Zusammenhang $R$ und $C$}{%
      R\cdot C=\tfrac{\varepsilon}{\kappa}}} &
  \beispielblock{Plattenkondens.\ $A$, $d$}{%
    \begin{aligned}
      C&=\varepsilon\tfrac{A}{d}\\[0.1em]
      G&=\kappa\tfrac{A}{d}\\[0.1em]
      R\cdot C&=\tfrac{\varepsilon}{\kappa}
    \end{aligned}} \\[0.3em]
\end{formelreihe}
\bereichende
